%%% Colour commands - delete for final draft 
% --- Editing Macros (Hierarchy of Importance) ---
% 1. FACT CHECK (Critical): Red commands immediate attention for potential inaccuracies.
\newcommand{\fc}[1]{\textcolor{red}{[FactC: #1]}} 
% 2. CITE (High): Orange highlights missing evidence or attribution.
\newcommand{\ct}[1]{\textcolor{orange}{[CITE: #1]}} 
% 3. APPENDIX (Medium): Blue indicates structural moves (shifting text out of the main body).
\newcommand{\app}[1]{\textcolor{blue}{[APP: #1]}} 
% 4. SMALL REWORD (Low): Teal is for minor stylistic edits, mirroring your friend's setup.
\newcommand{\sr}[1]{\textcolor{teal}{[Sreword: #1]}}
% 5. FIGURE (Low): Purple is for figure-related notes (refer to figure, create figure, etc.).
\newcommand{\fig}[1]{\textcolor{purple}{[FIG: #1]}}

%=======================================================================================
\chapter{Methodology and Experimental Design}
\label{ch:methodology}

This chapter outlines how the raw data was utilised to answer the three research questions. It is organised by method, first the the experimental framework, then common design decisions and data leakage control, the external environmental feature set construction, the top-height prediction models, the two attribution analyses, evaluation design, and finally assumptions, limitations and reproducibility. 

%------------------------------------------------
\section{Research questions and experimental framework}
\textbf{RQ1} test which models accurately predict top height and under what input and cohort year. \textbf{RQ2} is an attribution task, asking which areas over or under predict top height relative to the the average CR equation formed from the entire dataset. \textbf{RQ3} fits each plot with its own asymptote from only the plots repeated observations, and makes comparisons on the growth trajectories of curves, and if they are associated with specific stand, management or external environmental variables. 


\begin{table}[!htbp]
\scriptsize
\renewcommand{\arraystretch}{1.2}
\begin{tabular}{p{0.8cm}p{1.4cm}p{3.3cm}p{4.2cm}p{2.9cm}}
\toprule
RQ & Unit & Target & Model families compared & Evaluation design \\
\midrule
RQ1 & Plot-survey row & 95th-percentile LiDAR top height & Chapman-Richards curve, average-by-age
lookup, linear/random forest/XGBoost, deep neural network, Chapman-Richards-guided neural network
variants & Spatial-block (primary), pooled spatial $k$-fold, temporal, plot-level \\
RQ2 & Plot & Mean residual from the single pooled Chapman--Richards curve, across a plot's survey
history & Elastic Net, XGBoost, two-stage mixed-effects model & Spatial-block \\
RQ3 & Plot & Difference between a plot's own fitted asymptote and its yield-class-implied
asymptote & Elastic Net, XGBoost, GNNWR & Spatial-block \\
\bottomrule
\end{tabular}
\caption{Research questions mapped to unit of analysis, target, model families, and evaluation
design.}
\label{tab:rq-framework}
\end{table}

%------------------------------------------------
\section{Common Experimental Design and Data Leakage control}

\paragraph*{Evaluation Framework and Data Splits}. 4survey is the main cohort for every research question, used for the spatial-generalisation splits; 6survey (\S\ref{sec:data-cohorts}) is used for the temporal split and run in parallel elsewhere as a secondary check wherever sample size permits. For RQ1, each row in the dataset represents a single plot in a given survey year. For RQ2 and RQ3, data are aggregated so each row represents a single plot (Table~\ref{tab:rq-framework}). Model performance is evaluated using four distinct data-splitting strategies:

\begin{itemize}
\setlength\itemsep{0pt}
\item \textbf{Plot-level}. Individual plots are split at random (60/20/20, train/validation/test, seed 42), following previous work \ct{source that uses plot level split}. However a held out plots nearest neighbour is typically a training plot meters away, so a model can succeed by picking on its location in Aberfoyle rather than learning from features that shape growth trajectory. Kept as a baseline reference for comparison.
 
\item \textbf{Spatial block} (primary). Addresses spatial data leakage by assigning entire forestry compartments to train/validation/test sets (60/20/20, row-count-balanced, seed 42). To eliminate edge effects, any training plot within 60\,m of a held-out plot is excluded (validation/test plots are always retained; Figure~\ref{fig:split-maps}, bottom row). Primary isolation is achieved at the compartment level, while the 60\,m buffer—based on typical plot spacing, reduces local spatial autocorrelation along boundaries without discarding excessive data. \citep{roberts2017Cross_C5_EVAL}. 

\item {\textbf{Spatial block, pooled $k$-fold}} (stronger confirmation of the single spatial-block split). Compartments were assigned to five folds balanced approximately by row count. Each fold was held out once as the test set, the following fold was used for validation, and the remaining folds were used for training. Thus, every compartment received one prediction from a model that was not trained on that compartment. Held-out predictions from all five folds were pooled before calculating the overall performance metrics, avoiding distortions caused by differences in fold size and the range of observed heights\cite{ploton2020Spatial_C5_EVAL}. 

\item \textbf{Temporal} (secondary, harder test). Whole survey years, not compartments, are assigned to train/validation/test: 4survey trains on 2008 and 2012, validates on 2021, tests on 2023; 6survey trains on 2002/2006/2008/2012, validates on 2021, tests on 2023. This asks whether a model trained on earlier growth still predicts a real future survey it never saw, a different, harder question than spatial generalisation (Figure~\ref{fig:temporal-split}), following on from previous work \cite{lynch2025Digital_C4_NN}.
\end{itemize}

Every model uses the same split assignment for a given (cohort, split type, seed) combination, so results are always compared on identical train/validation/test membership. Compartment/block/sub-compartment identifier columns are never included in any model's feature list, enforced automatically. Shared data partitions and matched architectures reduce confounding when comparing model families under the same split e.g. the DNN and the Chapman-Richards-PINN. - though residual differences can still reflect optimisation behaviour or stochastic variation rather than the model design alone. %TODO: CHECK IF PINN ARCHTEICTURE IS SAME AS DNN. 


% FIGURE PLACEHOLDER (plot-level vs. spatial-block split, both cohorts) -- not finalised, likely
% candidate: plot_level_split map vs. spatial_block_split + buffer map, 4survey and 6survey.
% Style TBD -- match whatever map convention is settled on for the rest of the dissertation
% (e.g. the data chapter's cohort coverage map) rather than deciding it here in isolation.
% THREE-PANEL HORIZONTAL FIGURE PLACEHOLDER - also the temporal bar chart too. 
\begin{figure}[!htbp]
\centering
\begin{minipage}[b]{0.32\textwidth}
  \centering
  \fbox{\parbox{\textwidth}{\centering\vspace{2.5cm}\textbf{Panel A}\\Plot-Level Split\vspace{2.5cm}}}
\end{minipage}\hfill
\begin{minipage}[b]{0.32\textwidth}
  \centering
  \fbox{\parbox{\textwidth}{\centering\vspace{2.5cm}\textbf{Panel B}\\Spatial-Block Split ($60\text{~m}$ Buffer)\vspace{2.5cm}}}
\end{minipage}\hfill
\begin{minipage}[b]{0.32\textwidth}
  \centering
  \fbox{\parbox{\textwidth}{\centering\vspace{2.5cm}\textbf{Panel C}\\Temporal Split Timeline\vspace{2.5cm}}}
\end{minipage}
\caption{Overview of data-splitting strategies across cohorts: (A) Unconstrained random plot-level split showing nearby spatial leakage; (B) Compartment-based spatial-block split with a $60\text{~m}$ boundary exclusion buffer; and (C) Sequential temporal split timeline across survey years.}
\label{fig:splits-combined}
\end{figure}

\paragraph*{Preprocessing and leakage prevention}. Continuous features are standardised to zero mean and unit variance using parameters fitted exclusively on the training partition to prevent scaling data leakage. Age is scaled independently using a separate scaler, allowing its training-split standard deviation to be extracted directly for the physics-informed loss function's unit conversion (\S\ref{sec:rq1-pinn-loss}??).( Unordered categorical soil variables are one-hot encoded to avoid imposing artificial numeric orders.  -- make this sentnece simpler an shosrter. )Any plot with missing values in required features is removed entirely across all survey years to ensure models never encounter incomplete temporal trajectories. For the terrain-conditioned RQ1 models this removes 39 plots (156 rows) from 4survey and none from 6survey. \app{add to appendix which cells these were or where they are located on the map}


%------------------------------------------------
\section{Environmental feature-set construction}
Candidate environmental and management variables are screened separately for RQ1 (predictive tiers), RQ2 (shared-curve attribution), and RQ3 (plot-level attribution). All three follow the same recipe below, differing only in candidate pool and
target column. This produces five tiered feature sets (Set1–Set5; \S\ref{sec:feature-final}) to evaluate whether expanding environmental inputs improves model performance or introduces unhelpful noise across different model architectures.

\begin{enumerate}
\setlength\itemsep{0pt}
\item \textbf{Coverage screening.} Some variables are missing so some plots have no value at all for that variable, so the whole plot is removed which shrinks the datasets useable rows before any feature set is chosen. This removes 430 plot, from the 4-surveys candidate pool, mostly due to missing data from SoilGrids (413 plots missing), and CEH Topographic Wetness Index (39 plots, mostly
  clustered in one compartment \app{add these missing variable to appendix), plot of them. with examples plot ids.} The 6-survey subset contains no missing values. 

\item \textbf{Deduplication.} Deterministic duplicates (variables that are closed-form functions of another candidate) are dropped for example, Near-exact empirical duplicates (Spearman $|\rho| \geq 0.95$) are reduced to one variable, preferring externally measured variables over derived ones because they required fewer additional assumptions. 

\item \textbf{Importance ranking.} Three feature-importance signals: univariate Spearman $|\rho|$, permutation importance, and drop-column ablation (the change in held-out $R^2$ when removing a variable), are computed against the target variable and combined by rank-averaging. Correlation ignores context, permutation is biased under correlated predictors, ablation depends on one model/split, so we average across all three balancing out each methods limitations. 

\item \textbf{Variance-inflation screening.} To control for multicollinearity, variables are iteratively dropped if their VIF exceeds 5.0. Baseline columns are exempt from removal to anchor all feature sets (Sets 1–5), but still inform VIF calculations for other features.\citep{obrien2007Caution_C5_EVAL}. A
semivariogram, fitted to test-partition residuals only, complements Moran's I with the distance
range over which spatial autocorrelation persists, rather than only whether it is present. %% fix this sentence on the semiveriogram 
\end{enumerate}
% (turn above into table to save space if needed)
%% set 1 - 5 table should be after we explain the rsq so it makes more sense why they are different: 

\begin{table}[!htbp]
\centering
\scriptsize
\renewcommand{\arraystretch}{1.2}
\begin{tabular}{p{1cm}p{5.3cm}p{7.2cm}}
\toprule
Set & Selection rule & RQ1 membership (worked example) \\
\midrule
Set1 & Baseline only. & 4 columns: CanopyCover, time\_since\_thinning,
time\_since\_thinning\_missing, recent\_thinning\_5yr. \\
Set2 & Baseline + the top 10 candidates by combined importance rank (Spearman + permutation
importance + drop-column ablation, averaged); a candidate that fails VIF screening is skipped and
replaced by the next-ranked candidate, so the count stays at 10. & 14 columns: Set1 +
gwa\_weibull\_k\_50m, dist\_to\_scpt\_boundary, elevation, tas\_mean, eastness, windward\_topex,
gwa\_weibull\_a\_10m, gwa\_weibull\_a\_50m, slope\_degrees, cpmt\_compactness\_ratio. \\
Set3 & Baseline + the upper half (by combined rank) of the terrain/wind category only -- ranked
within that category, not against every other candidate; VIF-screened. & 15 columns: Set1 +
gwa\_weibull\_k\_50m, elevation, eastness, windward\_topex, gwa\_weibull\_a\_10m,
gwa\_weibull\_a\_50m, slope\_degrees, gwa\_weibull\_k\_10m, solar\_radiation\_index, ceh\_twi,
gwa\_wind\_speed\_10m. \\
Set4 & Set3 + the upper half (by combined rank) of every other category -- climate, soil/site,
spatial-position/edge-effects; VIF-screened. & 18 columns: Set3 + dist\_to\_scpt\_boundary,
tas\_mean, chelsa\_gdd5\_degc, cpmt\_compactness\_ratio. \\
Set5 & Baseline + every deduplicated candidate, with no importance-rank filter at all --
essentially the full candidate pool, cleaned only of duplicates and VIF-redundant members. & 30
columns: baseline + 26 further candidates spanning every category (e.g. slope\_degrees,
soilgrids\_ph, chelsa\_gdd5\_degc, dist\_to\_road, whcl); full list in
Appendix~\ref{app:feature-sets}. \\
\bottomrule
\end{tabular}
\caption{The five feature sets created by applying the screening steps above, here is the worked example for RQ1.
RQ2 and RQ3 apply the identical steps over their own candidate pools (different available variables,
different target column), therefore different final membership at each Set -- full membership
for all three research questions is given in Appendix~\ref{app:feature-sets}.}

\caption{The five feature sets generated by the screening protocol, using RQ1 as an example. While RQ2 and RQ3 follow the identical protocol, differing candidate pools and target variable for the models lead to different set membership (full lists in Appendix~\ref{app:feature-sets}).}
\label{tab:feature-sets-composition-rq1}
\end{table}

\begin{table}[!htbp]
\centering
\scriptsize
\renewcommand{\arraystretch}{1.2}
\begin{tabular}{p{2.3cm}p{4.2cm}p{1.6cm}p{2cm}p{3.4cm}}
\toprule
Feature set family & Feeds & Model comparison & Attribution & Purpose \\
\midrule
RQ1 & 30 & DNN and Chapman--Richards-guided network variants (environment-conditioned) & Primary &
No & DNN-vs-guided-network comparison under progressively richer terrain/wind/climate input; no
coefficient-based model reads these lists. \\
RQ2 & 28 & Elastic Net, XGBoost, two-stage mixed-effects model & Secondary & Primary & Which
conditions align with persistent departure from the shared curve. \\
RQ3 & 31 & Elastic Net, XGBoost, GNNWR & Secondary & Primary & Which conditions align with
plot-specific deviation; GNNWR additionally tests whether that association is spatially
varying. \\
\bottomrule
\end{tabular}
\caption{Purpose of the feature sets, each research question uses the identical Set1--Set5 screening protocol, with "Set5 size" reflecting the total feature count after deduplication and VIF screening}
\label{tab:feature-set-purpose}
\end{table}

% FIGURE PLACEHOLDER (pipeline overview -- simplified, high level only)
% Caption: "From plots to results: cohorts feed spatial/temporal partitions and a shared
% screening process producing Set1-5 per research question. Three branches follow: RQ1's
% predictive models, RQ2's shared-curve residual attribution, RQ3's plot-specific asymptote
% attribution, each with its own evaluation output."
%
% Confirmed evaluation output per branch (do not claim more than this without checking the
% notebooks under notebooks/spatial_analysis/, which were not inspected for this pass): RQ1
% produces per-fold and pooled predictive metrics under each split type; RQ2 produces a
% fixed-effects coefficient table plus held-out residual predictions; RQ3 produces per-fold
% predictive metrics, a per-plot coefficient table (GNNWR), and a residual spatial-autocorrelation
% diagnostic. No rendered coefficient/residual map was confirmed anywhere in the modelling
% pipeline -- do not draw one as an output box unless separately verified.
%
% Mermaid sketch (simplified, for redrawing): -- improve graphically.  also make sure its correct. 
% ```mermaid
% flowchart TD
%     A["Cohorts: 4survey / 6survey"] --> B["Spatial and temporal<br/>partitions"]
%     A --> C["Candidate environmental<br/>and management variables"]
%     C --> D["RQ-specific feature<br/>screening: Set1-5"]
%     B --> E1
%     B --> E2
%     B --> E3
%     D --> E1["RQ1: top-height<br/>prediction models"]
%     D --> E2["RQ2: shared-curve<br/>residual attribution"]
%     D --> E3["RQ3: plot-specific<br/>asymptote attribution"]
%     E1 --> F1["Predictive metrics,<br/>per-fold and pooled"]
%     E2 --> F2["Fixed-effects coefficient<br/>table + held-out residuals"]
%     E3 --> F3["Predictive metrics +<br/>per-plot coefficient table"]
% ```
\begin{figure}[!htbp]
\centering
\fbox{\parbox{0.9\textwidth}{\centering\vspace{4cm}FIGURE PLACEHOLDER --- pipeline overview, see
comment above source for exact content\vspace{4cm}}}
\caption{Overview of the full pipeline from raw data to the three research questions' results.}
\label{fig:pipeline-overview}
\end{figure}

%------------------------------------------------
\section{RQ1: Top-height prediction models}

RQ1 evaluates top-height models across four dimensions: baseline predictive accuracy, the use of supplementary environmental data (which can help stakeholders learn more where to invest into further high-quality data acquisition to improve modelling and planning), the benefit of Chapman-Richards physics guidance over unguided networks, and performance generalisation across spatial and temporal splits.

\textbf{Chapman-Richards (CR) curve.} A three-parameter curve
\begin{equation}
H(a) = y_{\max}\left(1-e^{-ka}\right)^{p}
\label{eq:cr}
\end{equation}
predicts top height $H(a)$ at age $a$ using asymptotic height $y_{\max}$, growth rate $k$, and shape parameter $p$. Curves are fitted strictly on each split's training partition via nonlinear least squares (five restarts, lowest sum of squared errors (SSE); Appendix~\ref{app:sse}). Under the pooled $k$-fold split, five separate curves are fitted per cohort (one per fold). The resulting $(y_{\max}, k, p)$ serve as fixed constants in the PINN loss (\S\ref{sec:rq1-pinn-loss}): each PINN model uses strictly its own fold-specific curve so that information from held-out validation or test compartments never leaks into the physics guidance loss.

\textbf{Average-by-age.} Mean observed top height is calculated for each integer age on the training partition, with unseen ages assigned the overall training mean. Mirroring conventional yield-table practice \citep{forestResearchForestYield}, this empirical baseline sets a performance threshold that complex models must exceed to provide predictive value beyond expected age-height averages.

Three standard supervised models are fit on all RQ1 feature sets (baseline and Sets 1–5) using identical splits. Linear regression fits one. 
\textbf{Linear regression.} A single linear relationship was fitted between the predictors and height. Because plots were at least 15 years old in 2008, the observed section of the age--height relationship may appear approximately linear despite the nonlinear full growth trajectory. Linear regression therefore provides a simple baseline, but may perform poorly near the limits of the observed age range. \app{but age distribution, the cut off for early ages (with annotation of unreliable lidar scans) and then an example champan richards curve, for the 4 point example and showing how it aligns with the linear portion of the curve}


\textbf{Random Forest.} Averages bootstrap-aggregated decision trees to capture nonlinear interactions without parametric growth assumptions. This provides a non-parametric inductive bias distinct from the Chapman--Richards and DNN/PINN families, testing whether top-height prediction requires smooth curve constraints or flexible ensemble averaging. (%TODO in background justify this baseline, cite commone sources e.g. lee2018MachineLearningPlotbased observatons doi     = {10.3390/f9050268}) 

\textbf{XGBoost} builds regression trees sequentially, with each tree correcting errors remaining from the preceding ensemble:
\begin{equation}
\hat{y}_i^{(K)} = \sum_{k=1}^{K}\eta f_k(\mathbf{x}_i),
\end{equation}
where $\mathbf{x}_i$ contains the predictors for observation $i$, $f_k$ is the $k$th regression tree, and $\eta$ is the learning rate. Unlike Random Forest, which averages independently fitted trees, this sequential approach can capture complex nonlinear relationships and provides a strong conventional machine-learning comparison for the neural models.

To ensure consistency across all research questions, a single hyperparameter configuration (max depth 4, learning rate 0.04, and up to 500 boosting rounds with early stopping) was applied to every XGBoost model across RQ1–RQ3, rather than tuning each feature set independently. This configuration was carried over from an earlier exploratory fit rather than derived from a systematic grid search. RQ3 used this setup from the outset, whereas RQ1 and RQ2 were initially run on library defaults and subsequently refit under this shared configuration for parity.\sr{rewrd this setnece, too many uses of word conftioufraion} A sensitivity analysis showing the impact of this adjustment on accuracy is provided in Appendix X, with full hyperparameter details in Appendix~\ref{app:hyperparameters}. \app{add to appendix)}

\subsection{Neural Network Models}
All DNN and CR-PINN models share an identical core architecture (Figure~\ref{fig:model-architecture}) to ensure that performance differences reflect the loss formulation rather than variations in network capacity. Instead of feeding directly into the base network, environmental covariates enter compact auxiliary sub-networks to adjust the physical curve parameters. This design follows the principles of conditional and variable-coefficient PINNs, where physical constraints or parameters vary with instance-level conditions \citep{kovacs2022Conditional_C4_NN,miao2023VCPINN_C4_NN}. Testing four alternative network sizes confirmed that this shared design performs best across both cohorts, with guided models proving the least sensitive to size changes (Appendix~\ref{app:architecture-screening}).

% FIGURE PLACEHOLDER (model architecture)
% Caption: "Shared DNN/guided-network architecture. The plain data loss (DNN, and the guided
% network's own data term) uses the main trunk's output directly. The guided network adds two
% consistency terms against the Chapman-Richards curve's analytical derivative: an instantaneous
% check via automatic differentiation (physics loss) and a finite-difference check across real
% survey pairs (trajectory loss). In the environment-conditioned variants, terrain/wind features
% feed only the small asymptote (and, in the rate-conditioned variant, rate) sub-networks --
% never the main trunk -- producing a per-plot adjustment to the CR curve's asymptote/rate that
% the physics and trajectory losses are then computed against."
% Content: a block diagram. Left: [Age, no-env features] -> main trunk (3x128, shared) ->
% [predicted height]. Below/beside: [Terrain/wind features] -> small asymptote sub-network (16
% units) -> [y_max adjustment, additive] and, dashed/optional for the rate-conditioned variant,
% -> rate sub-network -> [k adjustment, multiplicative, log-space]. A separate branch shows
% [predicted height] and [age] feeding into an autograd box labelled dH/da, compared against a
% box labelled "CR analytical derivative (frozen y_max, k, p)" to produce the physics loss; a
% second small diagram shows two real survey observations of one plot feeding two forward
% passes, differenced and compared to the same CR derivative at mid-age, to produce the
% trajectory loss. These survey pairs are pre-built once from real consecutive-survey rows before
% training starts (not synthesised per batch) -- draw them as a small fixed table/lookup feeding
% the diagram, not as an on-the-fly sampling step.

\begin{figure}[!htbp]
\centering
\fbox{\parbox{0.9\textwidth}{\centering\vspace{4cm}FIGURE PLACEHOLDER --- Model Architecture Block Diagram\vspace{4cm}}}
\caption{Shared architecture and loss formulation for the standard and Chapman--Richards-guided networks. All variants share a 3-layer predictive trunk ($3 \times 128$). In environment-conditioned models, terrain and wind features feed only auxiliary sub-networks (16 units) to adjust the asymptotic height ($y_{\max}$) and growth rate ($k$). The guided variants augment the standard data loss with: (1) an autograd-derived \emph{physics loss} ($\mathrm{d}H/\mathrm{d}a$) evaluated against the analytical derivative, and (2) a \emph{trajectory loss} computed over consecutive empirical survey pairs.}
\label{fig:model-architecture}
\end{figure}

The \textbf{Deep neural network} uses three hidden layers of 128 units each, leaky-ReLU activation, age concatenated with every other
input feature. The DNN's loss is plain mean squared error,
\begin{equation}
\mathcal{L}_{\text{data}} = \frac{1}{N}\sum_{i=1}^{N}\left(\hat H_i - H_i\right)^2,
\label{eq:data-loss}
\end{equation}
where $\hat H_i$ is the network's predicted height for row $i$ and $H_i$ is the observed height,
plus an L1 weight penalty ($\lambda_{L1}=10^{-5}$). The environment-conditioned DNN is
architecturally identical to the no-environment DNN; its terrain/wind block is simply concatenated
into the same input, widening that vector rather than adding a second network. This is so that interactions between age and stand variables can be learnt without imposing a separate structure. 


The \textbf{Chapman-Richards-guided neural network (CR-PINN)} models differ from the baseline DNN only by their loss function. In the baseline no-environment variant, the CR parameters $(y_{\max}, k, p)$ are treated as fixed cohort constants (Equation~\ref{eq:cr}) rather than learned parameters. The model optimises a three-term loss informed by the analytical growth derivative $H'(a) = y_{\max} p k e^{-ka}(1 - e^{-ka})^{p-1}$:

\begin{subequations}
\label{eq:cr-pinn-losses}
\begin{align}
\mathcal{L}_{\text{phys}} &= \frac{1}{N}\sum_{i=1}^{N}\left(\frac{\partial \hat{H}_i}{\partial a_i} - H'(a_i)\right)^2, \label{eq:physics-loss} \\
\mathcal{L}_{\text{traj}} &= \frac{1}{M}\sum_{j=1}^{M}\left(\frac{\hat{H}_{j,\text{later}} - \hat{H}_{j,\text{earlier}}}{\Delta a_j} - H'(\bar{a}_j)\right)^2, \label{eq:trajectory-loss} \\
\mathcal{L} &= \mathcal{L}_{\text{data}} + \lambda_{\text{phys}}\mathcal{L}_{\text{phys}} + \lambda_{\text{traj}}\mathcal{L}_{\text{traj}} + \lambda_{L1}\sum_{\theta}|\theta|. \label{eq:total-loss}
\end{align}
\end{subequations}

Here, the \emph{physics loss} \eqref{eq:physics-loss} checks the model's instantaneous growth rate via automatic differentiation at single points (rescaled by the training-split standard deviation to match real-age units). The \emph{trajectory loss} \eqref{eq:trajectory-loss} checks the predicted height gain across $M$ consecutive real survey visits from the same plot, where $\Delta a_j$ is the time elapsed and $\bar{a}_j$ is the mid-point age (drawn strictly from training rows to avoid leakage). The total objective \eqref{eq:total-loss} combines data loss, both consistency penalties (default $\lambda_{\text{phys}}=\lambda_{\text{traj}}=1.0$; ablation in \S\ref{sec:rq1-ablation}), and $L_1$ regularisation.

For the \textbf{Environment-conditioned growth} variants, terrain and wind features $\mathbf{x}^{\text{terrain}}_i$ feed into compact sub-networks (16 units) to produce plot-specific parameter adjustments, leaving the shape parameter $p$ fixed as a global constant:

\begin{subequations}
\label{eq:env-adjustments}
\begin{align}
y_{\max}^{(i)} &= y_{\max} + f_{y}\!\left(\mathbf{x}^{\text{terrain}}_i\right), \label{eq:ymax-adjust} \\
k^{(i)} &= k \cdot \exp\!\left(f_{k}\!\left(\mathbf{x}^{\text{terrain}}_i\right)\right). \label{eq:k-adjust}
\end{align}
\end{subequations}

Here, $f_y$ applies an additive adjustment to the asymptotic ceiling $y_{\max}$ \eqref{eq:ymax-adjust} following site-index modelling principles \citep{socha2023Higher_C1_SS}, while the rate-conditioned variant adds $f_k$ to apply a multiplicative adjustment to the growth rate $k$ \eqref{eq:k-adjust}. Both sub-networks receive only environmental features and are initialised near zero, ensuring training begins from the validated pooled curve before learning local adjustments.

\paragraph*{Hyperparameter sweeps and ablation.} To evaluate the true impact of physical guidance, the penalty weights are varied across $\lambda_{\text{phys}} = \lambda_{\text{traj}} \in \{0, 1, 2\}$ and compared against the baseline DNN of identical architecture. The symmetric default ($\lambda = 1.0$) from \eqref{eq:total-loss} was pre-specified before testing instead of being selected after to avoid circular evaluation. Consequently, all guided models in the primary comparison retain this fixed default (full experimental settings and results are detailed in the Results chapter).


\begin{table}[!htbp]
\scriptsize
\renewcommand{\arraystretch}{1.2}
\begin{tabular}{p{3.2cm}p{3cm}p{2.6cm}p{4.5cm}}
\toprule
Model & Environmental input & Loss terms & Notes \\
\midrule
CR (baseline) & None & N/A (curve fit) & Pooled per cohort, refit per fold for pooled $k$-fold. \\
Average-by-age & None & N/A (lookup) & Same training rows as CR. \\
Linear / random forest / XGBoost & No-environment feature set & Model-specific & XGBoost uses a
fixed configuration, not a validation-selected grid (Appendix~\ref{app:hyperparameters}). \\
DNN (no environment) & None & Data (Eq.~\ref{eq:data-loss}) + L1 & Shared architecture with every
guided variant. \\
DNN (environment-conditioned) & Terrain/wind, concatenated into main input & Data + L1 & Same
network class, wider input only. \\
Guided network (no environment) & None & Data + physics + trajectory + L1 & $y_{\max},k,p$
pooled, frozen. \\
Guided network (terrain-conditioned asymptote) & Terrain/wind, asymptote sub-network only & Data +
physics + trajectory + L1 & Eq.~\ref{eq:ymax-adjust}; $k,p$ pooled. \\
Guided network (terrain-conditioned asymptote and rate) & Terrain/wind, asymptote and rate
sub-networks & Data + physics + trajectory + L1 & Eqs.~\ref{eq:ymax-adjust}--\ref{eq:k-adjust};
$p$ pooled. \\
\bottomrule
\end{tabular}
\caption{Model family comparison for top-height prediction (RQ1). Full hyperparameters:
Appendix~\ref{app:hyperparameters}.}
\label{tab:rq1-models}
\end{table}

%maybe we paraphrase the metrics we will look at here then refer to the later evlaution metrics portion of what these are 

\section{RQ2a: does environment shrink RQ1's own departure from the curve}
RQ2a reuses CR-parameter predictions from RQ1 alongside the fold-matched Chapman--Richards baseline without fitting new models. Deviations from this baseline curve mark plots that diverge from the population average. We test whether environmental covariates reduce these departures, particularly for the most severely over- or under-predicted plots, rather than providing a uniform global improvement.

\paragraph*{Residual reduction}
For each test observation, predictions from the fold-matched CR baseline and the environment-conditioned models (DNN, PINN, and PINN$_k$) are matched by plot and survey year. The residual reduction is computed from the error magnitudes:
\begin{equation}
\text{reduction}_i = \left|r_i^{\text{CR}}\right| - \left|r_i^{\text{model}}\right|,
\label{eq:residual-reduction}
\end{equation}
where a positive value indicates the model predicts closer to the true height than the baseline curve. Evaluating the residual's sign ($r_i = y_i - \hat{y}_i$) alongside \eqref{eq:residual-reduction} identifies whether the model corrects directional over- or under-prediction by the shared curve, rather than merely shrinking error magnitude. 

\paragraph*{Quartile decomposition and coverage.} To separate general model improvements from specific environmental corrections, test observations are partitioned within each fold into equal quartiles by baseline error magnitude $|r^{\text{CR}}|$ (from Q1: smallest error, to Q4: largest). Because this split is purely rank-based, each quartile contains exactly 25\% of the test data; this evaluates performance across standard subgroups rather than focusing on a few extreme outliers. Instead of a single pooled mean, we can separate if environmental conditioning removes a small amount of error across all plots uniformly \textit{or} if conditioning can specifically reduce magnitude of errors on where the shared curve fits worse (Q4), capturing site-specific departures. Mean reductions are evaluated per quartile across all three environment-conditioned models, feature sets (Sets~2--4), and cohorts under spatial cross-validation, with the top-performing configuration verified across multiple random seeds for robustness. %try and define cohort as refring to 4 or 6 in the first place its mentioned in the diss. 

\section*{RQ2b: persistent departure from a shared growth curve.}
This asks which environmental/management conditions align with a plot's persistent position above or below the single shared curve, using only the 4survey cohort (6survey's 47 compartments were judged too few for a reliable spatial-block attribution split).

\textbf{Target construction.}
For each survey row $i$ in a plot's history, a residual is computed against the pooled baseline curve, and the attribution target $\overline{r}_{\text{plot}}$ is defined as the plot's mean residual across all $T$ observed surveys:

\begin{subequations}
\label{eq:target-construction}
\begin{align}
r_i &= H_i - H(a_i), \label{eq:cr-residual} & 
\overline{r}_{\text{plot}} &= \frac{1}{T}\sum_{t=1}^{T} r_{i,t}. \label{eq:mean-cr-residual}
\end{align}
\end{subequations}

Age is excluded from all candidate feature sets because it is the baseline curve's sole predictor in \eqref{eq:cr-residual}, which would introduce circularity into the target in \eqref{eq:mean-cr-residual}.

\paragraph*{Attribution models.}
Three models are fitted on $\overline{r}_{\text{plot}}$ across feature sets (Sets~2--4) under spatial-block cross-validation. Set~1 (baseline only) serves as a targeted ablation to test whether canopy cover signals reflect genuine environmental drivers or prior growth. 

\begin{enumerate*}[label=(\roman*)]
\item \textbf{Elastic Net} standardises continuous features, one-hot encodes soil classes, and tunes regularisation parameters ($\ell_1/\ell_2$ ratio and penalty strength) via 5-fold CV on the training partition.
\item \textbf{XGBoost} a separate model from RQ1's, fitted on RQ2's own target
  $\overline{r}_{\text{plot}}$ and RQ2's own candidate features rather than reused from RQ1;
  what carries over is only the hyperparameter recipe, not the fitted model itself (500 boosting
  rounds, maximum tree depth 4, learning rate 0.04, early stopping against a held-out validation fold.) %add to limiations on useing the same hyper parmetrs. for different rq? 
\item \textbf{Two-stage mixed-effects model.} Uses a two-stage approximation to a full non-linear mixed-effects model \citep{socha2016Assessment_C2_CR}. Stage~1 isolates persistent plot-level departures ($\overline{r}_{\text{plot}}$) from the frozen baseline curve, and Stage~2 fits a linear mixed-effects model against candidate environmental and management features:
\begin{equation}
\overline{r}_{\text{plot}} = \boldsymbol{\beta}^{\top}\mathbf{x}_{\text{plot}} + u_{\text{cpmt}} + \varepsilon_{\text{plot}},
\label{eq:mixed-effects}
\end{equation}
where $\boldsymbol{\beta}$ represents fixed-effect coefficients on standardised features $\mathbf{x}_{\text{plot}}$, $u_{\text{cpmt}} \sim \mathcal{N}(0, \sigma_u^2)$ is a compartment-level random intercept accounting for spatial clustering, and $\varepsilon_{\text{plot}}$ is the residual error. The model uses random intercepts only; random slopes were omitted because compartment sizes vary widely (from single plots to over a thousand) and no varying-effect hypotheses were posed. Fitted on the training partition just for attribution rather than held-out prediction, coefficients are estimated via restricted maximum likelihood (REML). Nested fixed-effect comparisons use maximum likelihood (ML), as REML likelihoods are not comparable across differing fixed-effects structures \citep{pinheiro2000Mixed_C5_EVAL}. Because $\overline{r}_{\text{plot}}$ is a pre-calculated residual, uncertainty from Stage~1 is not propagated; coefficients are interpreted strictly as statistical attribution of curve departures rather than causal mechanisms.
\end{enumerate*}
%possibly make into a table to save space. 

Feature importance and SHAP values provide post-hoc interpretability (detailed in \S\ref{sec:eval-diagnostics}). While gain-based importance is extracted directly during tree construction, SHAP values are computed post-fitting across the full dataset using the final model checkpoint. We compare the global ranking of environmental features against the linear mixed-effects coefficients to evaluate whether non-linear tree models and parametric baselines identify consistent environmental drivers, with full evaluation metrics detailed in \S\ref{sec:eval-diagnostics}. -- %change this sentence paragraph to end this section - refer to the actual evaluation metrics use but refer to the next section where these are explained. 

\section{RQ3: Plot-specific growth-curve deviation}
This tests if a plot's fitted growth ceiling deviates from its static yield-class expectation, and whether local environmental or management conditions explain this deviation. Deviation is estimated specifically in the asymptotic ceiling while keeping the curve's shape parameters fixed, for three reasons. First, because plots contain few survey observations, fixing the shape parameters allows the asymptote to be estimated via a stable, closed-form linear fit rather than an unstable non-linear optimisation. Second, the yield-class comparison directly concerns maximum height potential, making the asymptote the exact parameter of interest. Third, this follows standard site-index modelling conventions where site quality shifts asymptotic potential while trajectory shape remains constant \citep{socha2021Regional_C2_CR}. 

By construction, this target evaluates asymptotic height potential alone: plots with differing growth rates or inflection timing that ultimately reach the same eventual ceiling are not treated as divergent. \sr{refer to limitaitons}

\paragraph*{Target construction.}
For each plot, the asymptote $\hat{y}_{\max}^{\text{plot}}$ is estimated from its own repeated height--age surveys via closed-form weighted least squares through the origin, fixing the curve's shape parameters to their Forestry Commission yield-class lookup values. The target is defined as the deviation from the static yield-class asymptote:
\begin{equation}
\Delta y_{\max}^{\text{plot}} = \hat y_{\max}^{\text{plot}} - y_{\max}^{\text{yldc}}.
\label{eq:local-ymax-diff}
\end{equation}
This metric is constructed strictly from height--age records and static tables, with no environmental variables involved. 

Plots exhibiting extreme $\Delta y_{\max}^{\text{plot}}$ deviations were audited prior to modelling against independent disturbance flags (clearfelling, measurement anomalies, ambiguous disturbance \app{add rules to appendix}) and a trajectory-instability cutoff (top 1\% within-plot height change). Because no checked plots triggered these exclusion criteria, all were retained as genuine empirical departures rather than data artifacts.\fc{check this is true}

\paragraph*{Global and spatially varying attribution models.}
Elastic Net and XGBoost are fitted on $\Delta y_{\max}^{\text{plot}}$ using Sets~2--4 following the same protocols as RQ2b \cite{zou2005Regularization_C4_NN, chen2016XGBoost_C4_NN}. To test whether environmental relationships vary across the landscape rather than following a single global trend, we additionally fit Geographically and Neurally Weighted Regression (GNNWR) \cite{du2020Geographically_C6_GWR}):
\begin{equation}
\hat{y}_i = \beta_0(s_i) + \sum_j \beta_j(s_i)\,x_{ij},
\label{eq:gnnwr}
\end{equation}
where $s_i$ is plot $i$'s spatial coordinate and each $\beta_j(s_i)$ is a spatially varying coefficient. 

GNNWR learns spatial relationships via a neural network. First, a global OLS baseline is fitted and frozen. For any plot, the spatial sub-network takes the vector of Euclidean distances to all training reference plots and jointly predicts multiplicative scaling weights for the global coefficients. Local parameters $\beta_j(s_i)$ are formed by multiplying each frozen global OLS slope by its plot-specific learned weight. Test plots are evaluated using their distances to this fixed training set without memory capping. \fig{refer to figure rq3-attribution}

XGBoost attribution is evaluated using gain importance and SHAP values
\citep{lundberg2017}. Unlike the full-sample analysis in RQ2b, SHAP values in
RQ3 are calculated out of fold across the five spatial-block splits and then
pooled. Each plot's attribution is therefore obtained from a model that was
not trained on its spatial fold. (\S\ref{sec:eval-diagnostics})

% FIGURE PLACEHOLDER (RQ3 target construction and attribution models, incl. GNNWR mechanism)
% Caption: "RQ3's asymptote-deviation target, fitted from each plot's own repeated observations
% against a static yield-class expectation, feeds three attribution models: global Elastic Net
% and XGBoost, and GNNWR, whose spatial weighting network re-weights a frozen global OLS fit
% using each plot's distances to a training reference set, producing plot-level local
% coefficients and predictions evaluated under the spatial-block split. RQ2's parallel but
% simpler pipeline (pooled residual -> Elastic Net/XGBoost/two-stage mixed-effects model) is
% described in prose and equations only above, not duplicated in a figure here,
% since it has no comparable internal mechanism to unpack."
%
% Verified against the actual GNNWR implementation (third-party `gnnwr` package) before drawing:
% - input to the spatial weighting network is a vector of Euclidean distances from the focal
%   plot to every plot in the reference set (MinMax-scaled), NOT raw coordinates and NOT
%   coordinate differences;
% - no k-nearest-neighbour or distance-threshold step exists -- "neighbourhood" is implicit,
%   learned by the network from the full distance vector to the reference set;
% - the network does NOT output final coefficients directly: it outputs one multiplicative
%   weight per predictor (+ intercept), which is multiplied against a separately pre-fitted,
%   frozen global OLS coefficient to give the local coefficient beta_j(s_i);
% - one shared network with a multi-output head produces every coefficient's weight jointly, not
%   one network per coefficient;
% - held-out/test plots reuse the training reference set for their distance vector, not a
%   bespoke neighbour set built from their own location;
% - this is the standard GNNWR (Du et al. 2020) formulation; the reference set CAN be capped at a
%   fixed row count for memory reasons, but the reported runs use the full, uncapped training
%   partition (confirmed: run with --reference-set-size 0) -- do not draw a capped reference set;
% - no geographical kernel or bandwidth parameter exists in the implementation -- do not draw one;
% - no rendered spatial coefficient map was found in the modelling pipeline, only a per-plot
%   coefficient table and a Moran's I residual-autocorrelation diagnostic -- label the coefficient
%   output as a "per-plot coefficient table", not a "coefficient map", unless a map-producing
%   script is confirmed separately.
%
% Mermaid sketch (for redrawing, not final layout):
% ```mermaid
% flowchart TD
%     A["Repeated heights, own plot only"] --> B["Closed-form per-plot<br/>asymptote fit"]
%     C["Yield-class asymptote<br/>(static lookup)"] --> D["RQ3 target: y_max diff"]
%     B --> D
%     F["Environmental/management<br/>predictors (Set2-4)"] --> G["Elastic Net"]
%     F --> H["XGBoost"]
%     F --> M["Global OLS fit<br/>(frozen coefficients)"]
%     D --> G
%     D --> H
%     D --> M
%     P["Plot location"] --> Q["Distance to every plot in<br/>reference set (full training partition)"]
%     Q --> R["Spatial weighting network<br/>(one shared trunk)"]
%     R --> S["Per-predictor weights w_j(s_i)"]
%     M --> T["Local coefficient<br/>beta_j(s_i) = w_j(s_i) x beta_j"]
%     S --> T
%     T --> U["Plot-level prediction"]
%     T --> V["Per-plot coefficient table"]
%     G --> W["Spatial-block holdout evaluation"]
%     H --> W
%     U --> W
% ```
%
% Content (prose fallback if not drawn from the mermaid sketch above): left-to-right, target
% construction (closed-form per-plot asymptote fit vs. static yield-class asymptote) feeding
% three model boxes -- Elastic Net, XGBoost, and an expanded GNNWR box showing the
% distance-to-reference-set input, spatial weighting network, per-predictor weights, and their
% multiplication against the frozen global OLS coefficients -- all three converging on a
% spatial-block holdout evaluation box. A small annotation: "association, not causation."

\begin{figure}[!htbp]
\centering
\fbox{\parbox{0.9\textwidth}{\centering\vspace{4cm}FIGURE PLACEHOLDER --- RQ3 target construction
and attribution models, see comment above source for exact content\vspace{4cm}}}
\caption{Plot-specific asymptote-deviation target construction (RQ3) and its three attribution
models, including the GNNWR spatial weighting mechanism.}
\label{fig:rq3-attribution}
\end{figure}

\section{Model evaluation, uncertainty and interpretation} \label{sec:eval-diagnostics}
This section states what is measured and why, not what was found (Results chapter).

\paragraph*{Predictive diagnostics (RQ1).}
Every RQ1 model is scored on residuals $e_i = H_i - \hat H_i$ over $N$ rows:
\begin{align*}
\text{MAE} &= \frac{1}{N}\sum_i |e_i|, &
\text{RMSE} &= \sqrt{\frac{1}{N}\sum_i e_i^2}, \\
R^2 &= 1 - \frac{\sum_i e_i^2}{\sum_i (H_i - \bar H)^2}, &
\text{Bias} &= \frac{1}{N}\sum_i e_i.
\end{align*}
$R^2$ is variance explained relative to a mean-only baseline (1 perfect, 0 matches the mean,
negative is worse than the mean on a held-out split); positive bias is systematic
under-prediction. Metrics are further broken down by five age bands to check whether an
aggregate score conceals age-specific weaknesses (the maturity gate is applied at plot level, so
individual rows below 30 can still appear; a poor $<30$ score reflects a small, non-representative
sample, not necessarily worse prediction). Pooled $R^2$ additionally carries a cluster-bootstrap
95\% confidence interval, resampling whole compartments rather than individual rows or plots. This
is RQ1 only; RQ2b and RQ3's own $R^2$ numbers do not currently carry one.


\paragraph*{Spatial autocorrelation (Moran's $I$).}
Tests whether nearby plots have more similar residuals than distant ones, evaluated against a null hypothesis of spatial randomness. Global Moran's $I$ for model residuals $e_i = y_i - \hat{y}_i$ is defined as:
\begin{equation}
    I = \frac{N}{W} \frac{\sum_{i=1}^{N} \sum_{j=1}^{N} w_{ij} (e_i - \bar{e})(e_j - \bar{e})}{\sum_{i=1}^{N} (e_i - \bar{e})^2}
    \label{eq:morans_i}
\end{equation}
where $N$ is the total number of plots, $\bar{e}$ is the mean residual, $w_{ij}$ is the spatial weight between plots $i$ and $j$, and $W = \sum_{i=1}^N \sum_{j=1}^N w_{ij}$ is the sum of all spatial weights (\app{specify the spatial weighting scheme used, e.g., inverse distance threshold or $k$-nearest neighbours}). Values near zero indicate spatial randomness, positive values reflect spatial residual clustering (positive autocorrelation), and negative values indicate spatial dispersion.

Moran's $I$ is used in two distinct ways across this project: \textbf{RQ1 (Post-fit Diagnostic):} Evaluated on raw-height residuals to diagnose whether unresolved spatial structure remains after baseline fitting (used purely diagnostically without feeding back into feature selection). \textbf{RQ3 (Category-Removal Check):} Used in ablation tests to assess whether removing an environmental category increases residual clustering, diagnosing whether that feature group captured authentic spatial signal (currently computed on Set4 across the 4-survey cohort).

Moran's $I$ has not yet been computed on RQ2b attribution residuals, nor on RQ3's main GNNWR, Elastic Net, or XGBoost residuals under the current VIF-screened tiers. %TODO in rsults and evaluation -- edit this comment. once its done. 
%ADD IN SEMIVIOGRAM AND WHERE THAT IS IN THE METHOD. (i put it under vif comment, maybe it goes here, not sure?) 

\paragraph*{Feature attribution (RQ2b, RQ3).}
Feature importance is evaluated using XGBoost gain and SHAP values \citep{lundberg2017unified}, which decompose individual predictions into signed, additive feature contributions analogous to linear model coefficients. Neither metric is used for feature selection. Their analytical roles differ across the two questions:
\begin{itemize}
  \item \textbf{RQ2b:} SHAP provides a non-linear global attribution view to assess convergence alongside Elastic Net and mixed-effects coefficients.
  \item \textbf{RQ3:} SHAP additionally enables local attribution for individual outlier plots, isolating specific drivers of extreme growth deviations.
\end{itemize}
In both questions, SHAP values are pooled across the five spatial folds to report mean feature contributions; fold-to-fold SHAP stability is not evaluated.

\paragraph*{Coefficient stability (RQ2b, RQ3).}
Elastic Net coefficients are reported as mean $\pm$ standard deviation across the five spatial folds in both RQ2b and RQ3 to confirm that parameter signs and magnitudes remain consistent across spatially distinct partitions rather than reflecting a single data split. Attribution models are evaluated across spatial folds without multi-seed sweeps.

\paragraph*{Robustness to stochastic and configuration choices:} 
\begin{itemize}
\setlength\itemsep{0pt}
\item \textbf{Training-seed reseed} (RQ1, RQ2a): the RQ1 winner is refit across 4 additional
  seeds; RQ2a reuses these same predictions rather than reseeding separately, so the two share one
  robustness result.
\item \textbf{Architecture-size sensitivity} (RQ1, all three guided-network variants): tests
  whether the fixed network size affects the headline comparison.
\item \textbf{Physics-weight ablation} (RQ1, PINN and PINN$_k$): tests whether the guided
  network's advantage, if any, depends on the physics/trajectory loss weighting rather than the
  mechanism.
\item \textbf{XGBoost hyperparameter sensitivity} (RQ1 baseline, RQ2b, RQ2a): RQ1's baseline and
  RQ2b's attribution model were both found to be running on library defaults rather than a
  deliberately chosen configuration; this check tests whether their headline results depend on
  that choice, by refitting all three under one consistent, non-default configuration. Outcome
  reported in the Results chapter, not here.
\item \textbf{Split-seed reseed} (RQ3, GNNWR, 6survey only): checks whether GNNWR's
  smaller-cohort result is stable across different fold assignments, given its wide per-fold
  variance there, discussed as a limitation later in this chapter.
\end{itemize}


\section{Methodological assumptions and limitations}

% Keep to 3-5 items. Draft points below, expand each into 2-4 sentences by hand.

\begin{itemize}
\setlength\itemsep{0pt}
\item \textbf{Spatial-block and buffer assumptions.} The 60\,m buffer is set from typical plot
  spacing, not from the measured spatial-autocorrelation range of the shared-curve residual;
  compartment-level holdout is the primary leakage control, the buffer a
  secondary one.
\item \textbf{Environmental scale mismatch.} Terrain/wind/soil variables are raster-derived and
  assigned to point-scale plots; the raster resolution does not necessarily match the scale at
  which these conditions act on a plot.
  % TODO: confirm raster resolution(s) and the assignment method (nearest-cell, bilinear, zonal
  % mean) against the actual extraction code before stating this as fact.
\item \textbf{Feature-selection leakage or split dependence.} Ranking and VIF screening, described
  earlier, determine Set1--Set5 before final evaluation; the open questions
  flagged there about which partition was used, and whether ranking was fold-specific, bear
  directly on how much confidence to place in the resulting sets.
\item \textbf{Uncertainty in plot-specific targets and their attribution models.} RQ3's per-plot
  asymptote target is fitted from as few as the cohort minimum of repeated
  observations per plot; plots with shorter survey histories contribute a noisier estimate of the
  same target as plots with longer histories, without that difference being explicitly weighted or
  flagged downstream. This compounds with 6survey's small compartment count: a multi-seed reseed
  of RQ3's GNNWR fit on 6survey shows its pooled $R^2$ sign-flipping across seeds for every feature
  set, centred near zero within the seed-to-seed spread. This is an underpowered result, not a
  reliable positive or negative one. 6survey's GNNWR fit should be read as inconclusive, not as
  evidence the spatially-varying-coefficient effect is genuinely weaker there than on 4survey.
\item \textbf{Association, not causation.} Both attribution avenues report statistical association
  between environmental/management conditions and departure from
  an expected growth curve; none of the designs support a causal claim.
\end{itemize}

\section{Reproducibility and implementation}
--- \sr{add many more implementations from throughout this methodlogy, understand the limiations and how this affects resutls, and why we still chose to do it (still good choice, time limits, maybe another choice was bad so we did the one we explained) etc }

All split assignments, network weight initialisation, and Elastic Net/XGBoost cross-validation
folds use a fixed seed (42) unless explicitly swept as part of a stated robustness check (e.g. the
RQ1 winner-model reseed check, described above). Every model reads its feature list
from a single generated manifest, so a reported set's exact column membership can always be
checked against that source rather than a hand-copied list.

Detailed hyperparameter grids and selected values (Appendix~\ref{app:hyperparameters}), full
Set1--Set5 column lists and per-variable screening diagnostics (Appendix~\ref{app:feature-sets}),
and secondary sensitivity settings not reported in the main text (dropout-rate/architecture-width
sweeps, the narrow-gap temporal split's own results, Appendix~\ref{app:sensitivity}) are kept
out of this chapter and pointed to at the relevant section above rather than repeated here.
 